\documentclass{../../plantilla/raysac-informe}

\producto{transbank-SDK}
\titulo{Informe Técnico de Arquitectura, Protocolo y Diseño}
\titulocorto{Informe Técnico}
\version{1.0.0}
\fecha{2026-11-22}
\clasificacion{Uso interno}
\autor{raysac}

\begin{document}

\portada

\begin{controlversiones}
  \cvfila{1.0}{2026-11-22}{raysac}{Versión inicial exhaustiva del informe técnico: arquitectura sans-IO, protocolo de tramas, modelo de datos, API pública, seguridad, pruebas y runbooks.}
\end{controlversiones}

\avisolegalestandar

\clearpage
\preliminares

\section{Resumen ejecutivo}

\subsection{Propósito y visión técnica del proyecto}
El proyecto \texttt{transbank-SDK} implementa un kit de desarrollo de software (SDK) no oficial en el
lenguaje de programación Rust para la integración directa, segura y de alto rendimiento con terminales
físicos de punto de venta (POS) operados por la empresa adquirente Transbank en Chile. La biblioteca
soporta los protocolos normativos de los dos modelos comerciales dominantes en el mercado chileno:
\textbf{POS Integrado} (terminales atendidos en caja) y \textbf{POS Autoservicio} (kioscos y máquinas
desatendidas).

El proyecto fue concebido para resolver la ausencia de un SDK moderno, seguro y tipado en Rust que
permita vincular software de escritorio y sistemas embebidos de comercio con el hardware de pago
mediante interfaces serie físicas (RS-232 y USB CDC-ACM) o puentes de red TCP. El diseño garantiza
inmunidad contra estados de fallo ambiguos, asegurando de forma determinista la protección del
comercio y del consumidor.

\subsection{Motivación dentro del ecosistema retail}
La motivación principal del proyecto surge de la necesidad de dotar a la aplicación de punto de venta
\texttt{rework-sistema-facturacion} (sistema de escritorio basado en el framework Tauri y Rust) de una
capacidad de cobro con tarjeta ágil y confiable. Siguiendo las directrices estratégicas del ecosistema
raysac, se determinó desarrollar esta solución como un módulo de infraestructura completamente genérico,
independiente y desacoplado, sin tipos propietarios de facturación, liberado bajo doble licencia
código abierto (\texttt{MIT OR Apache-2.0}) para beneficio de la comunidad tecnológica.

\subsection{Principios rectores del diseño}
La arquitectura de \texttt{transbank-SDK} se fundamenta en cuatro principios normativos e inviolables:
\begin{enumerate}
  \item \textbf{Arquitectura Sans-IO primero (ADR-01):} toda la lógica de encuadre, análisis de tramas,
    cálculo de redundancia longitudinal (LRC) y formateo de mensajes reside en módulos puros, libres de
    llamadas al sistema operativo, entrada/salida o reloj. Esto habilita pruebas deterministas basadas en
    propiedades (\texttt{proptest}) y permite que el simulador reutilice el código exacto de producción.
  \item \textbf{Regla de oro: ``Nunca Doble Cobro'' (ADR-04, LINK-03):} una vez que un comando de venta o
    anulación ha sido acusado con un byte ACK por el POS, el SDK prohíbe terminantemente reenviar la trama
    ante un timeout o interrupción de canal. El resultado se declara formalmente como desconocido y se
    fuerza un protocolo de reconciliación mediante la consulta de última venta (\texttt{last\_sale}).
  \item \textbf{El rechazo financiero no es un error (ADR-03):} una tarjeta rechazada por falta de fondos,
    clave errónea o bloqueo es un resultado normal de negocio. La API devuelve \texttt{Ok(SaleResponse)}
    donde el método \texttt{is\_approved()} evalúa a \texttt{false}. El tipo \texttt{Err} se reserva
    exclusivamente para anomalías de comunicación o violaciones del protocolo serie.
  \item \textbf{Robustez estricta sin pánicos ni código inseguro (SAFE-01, AGENTS.md §6):} el workspace
    prohíbe de forma absoluta el uso de \texttt{unsafe} (\texttt{\#![forbid(unsafe\_code)]}) y el uso de
    métodos que provoquen pánico (\texttt{unwrap}, \texttt{expect}, \texttt{panic!}) en código de producción,
    utilizando buffers acotados para impedir desbordamientos de memoria.
\end{enumerate}

\subsection{Resumen de métricas del repositorio}
La Tabla~\ref{tab:metricas-repo} resume el estado técnico cuantitativo verificado en el repositorio.

\begin{longtable}{@{}p{6.5cm}p{4.5cm}p{3.5cm}@{}}
\caption{Métricas técnicas verificadas del workspace transbank-SDK.}\label{tab:metricas-repo}\\
\toprule
\textbf{Métrica / Componente} & \textbf{Valor verificado} & \textbf{Referencia} \\
\midrule
\endfirsthead
\toprule
\textbf{Métrica / Componente} & \textbf{Valor verificado} & \textbf{Referencia} \\
\midrule
\endhead
\bottomrule
\endlastfoot
Crates en el workspace & 3 (\texttt{pos}, \texttt{sim}, \texttt{cli}) & \texttt{Cargo.toml} \\
Edición de Rust / MSRV & Edición 2024 / MSRV 1.85 & \texttt{Cargo.toml} \\
Total de pruebas automatizadas & 401 pruebas (unitarias, integración, doc) & \texttt{cargo test} \\
Mutantes analizados (\texttt{cargo-mutants}) & 654 mutantes (0 escapes) & \texttt{TEST-01} \\
Lints de compilación y Clippy & Deny estricto sin warnings & \texttt{[workspace.lints]} \\
Licencia de distribución & Dual MIT / Apache-2.0 & \texttt{LICENSE-*} \\
\end{longtable}

\clearpage
\section{Definición del problema, actores y alcance}

\subsection{Contexto de pagos presenciales en Chile}
En el mercado minorista chileno, la captura presencial de tarjetas bancarias (crédito, débito Redcompra
y prepago) se efectúa mayoritariamente mediante terminales físicos provistos por Transbank. Existen dos
modalidades operativas:
\begin{itemize}
  \item \textbf{POS Autónomo (Tradicional):} el cajero digita manualmente el monto en el teclado del POS.
    Esta modalidad introduce errores humanos de digitación, descuadraturas de caja y lentitud en la atención.
  \item \textbf{POS Integrado / Autoservicio:} el software de punto de venta (POS host) envía automáticamente
    el monto exacto y el número de boleta al terminal a través de un cable USB o RS-232, eliminando la
    digitación manual y garantizando la cuadratura automática de los pagos.
\end{itemize}

\subsection{Actores del sistema y responsabilidades}
El sistema interactúa con seis entidades lógicas y humanas:
\begin{itemize}
  \item \textbf{Comercio:} entidad comercial titular de un contrato de afiliación con Transbank, identificada
    por un código de comercio (\texttt{commerce\_code}) de hasta 12 dígitos numéricos.
  \item \textbf{Host (Computador de Caja):} equipo informático que ejecuta el software de punto de venta y
    utiliza este SDK como biblioteca de enlace para solicitar cobros.
  \item \textbf{Terminal POS:} dispositivo de hardware seguro (Verifone VX520, Ingenico Desk 3500) certificado
    por Transbank. Lee el chip/banda/NFC, solicita el PIN y se comunica directamente con el switch adquirente.
  \item \textbf{Cliente (Pagador):} titular de la tarjeta que autoriza la transacción mediante su clave secreta.
  \item \textbf{Switch Transbank:} red adquirente bancaria que aprueba o rechaza los cargos financieros.
  \item \textbf{Desarrollador / Integrador:} ingeniero que programa sobre la API pública del SDK.
\end{itemize}

\subsection{Modalidades de terminal: Integrado vs Autoservicio}
La Tabla~\ref{tab:terminal-comparacion} compara las características técnicas y operativas de ambos modelos.

\begin{longtable}{@{}p{3.2cm}p{5.5cm}p{5.5cm}@{}}
\caption{Comparación técnica entre POS Integrado y POS Autoservicio.}\label{tab:terminal-comparacion}\\
\toprule
\textbf{Característica} & \textbf{POS Integrado} & \textbf{POS Autoservicio} \\
\midrule
\endfirsthead
\toprule
\textbf{Característica} & \textbf{POS Integrado} & \textbf{POS Autoservicio} \\
\midrule
\endhead
\bottomrule
\endlastfoot
Velocidad serie por defecto & 115.200 baudios & 19.200 baudios \\
Entorno operativo & Caja atendida por cajero & Kioscos, tótems y máquinas desatendidas \\
Impresión de vouchers & Impreso directamente por el POS & Impreso en el POS o devuelto al host \\
Procedimiento de arranque & No requiere inicialización previa & Requiere comando de inicialización (1000) \\
Comando de venta & Código \texttt{0200} & Código \texttt{1000} \\
Comando de anulación & Por número de operación (\texttt{1200}) & Anula solo la última venta (\texttt{1020}) \\
Totales y Detalle del día & Soportados (\texttt{0700}, \texttt{0260}) & No soportados por el protocolo \\
\end{longtable}

\subsection{Alcance técnico y exclusiones expresas}
Para mantener un enfoque de seguridad y especialización rigurosa, se establecen las siguientes delimitaciones:
\begin{itemize}
  \item \textbf{No es Webpay:} el SDK no interactúa con servicios web, REST APIs ni soluciones de comercio
    electrónico (Webpay Plus, Oneclick, Transbank REST).
  \item \textbf{Fuera del alcance PCI DSS (SEC-01):} el SDK nunca solicita, recibe ni almacena el número de
    tarjeta completo (PAN), el código de verificación (CVV) ni la clave personal (PIN). Toda la captura de
    datos sensibles ocurre dentro del módulo de seguridad física del terminal (PIN pad).
  \item \textbf{Sin lógica tributaria:} el SDK no genera boletas electrónicas ni documentos tributarios (DTE)
    ni se conecta con el Servicio de Impuestos Internos (SII).
  \item \textbf{Sin certificación oficial de marca (LEG-03):} este software es un desarrollo independiente
    no oficial. No reemplaza los procesos de homologación que Transbank pueda requerir a comercios específicos.
\end{itemize}

\clearpage
\section{Requisitos del sistema}

\subsection{Requisitos funcionales}
La Tabla~\ref{tab:requisitos-funcionales} enumera los requisitos funcionales implementados en el SDK.

\begin{longtable}{@{}p{1.4cm}p{7.2cm}p{2.2cm}p{3.4cm}@{}}
\caption{Requisitos funcionales del sistema transbank-SDK.}\label{tab:requisitos-funcionales}\\
\toprule
\textbf{ID} & \textbf{Descripción funcional} & \textbf{Prioridad} & \textbf{Estado} \\
\midrule
\endfirsthead
\toprule
\textbf{ID} & \textbf{Descripción funcional} & \textbf{Prioridad} & \textbf{Estado} \\
\midrule
\endhead
\bottomrule
\endlastfoot
RF-01 & Enmarcado y desenmarcado de tramas STX/ETX con checksum LRC. & P0 & Implementado (PROTO-01) \\
RF-02 & Parsing de mensajes y respuestas del POS Integrado. & P0 & Implementado (PROTO-02) \\
RF-03 & Control de enlace con ACK, NAK y reintentos automáticos acotados. & P0 & Implementado (LINK-01) \\
RF-04 & Manejo de temporizadores configurables de respuesta y acuse de recibo. & P0 & Implementado (LINK-01) \\
RF-05 & Filtrado y notificación de tramas intermedias de estado (0900). & P0 & Implementado (LINK-01) \\
RF-06 & Cliente de alto nivel para POS Integrado con métodos síncronos tipados. & P0 & Implementado (CLIENT-01) \\
RF-07 & Comunicación por puerto serie físico con autodetección de terminales. & P1 & Implementado (SERIAL-01) \\
RF-08 & Comunicación por socket TCP para puentes de red y pruebas locales. & P1 & Implementado (SERIAL-01) \\
RF-09 & Interfaz de línea de comandos (\texttt{tbk-pos}) con salida humana y JSON. & P1 & Implementado (CLI-01) \\
RF-10 & Simulador de terminal POS sans-IO con inyección de fallas y servidor TCP. & P0 & Implementado (SIM-01) \\
RF-11 & Soporte de operaciones completas para terminales POS Autoservicio. & P2 & Implementado (AUTO-01) \\
RF-12 & Cliente asíncrono con runtime Tokio (\texttt{AsyncPosIntegrado}). & P2 & Implementado (ASYNC-01) \\
RF-13 & Consulta de última venta (\texttt{last\_sale}) para reconciliación post-falla. & P0 & Implementado (CLIENT-01) \\
RF-14 & Anulación tipada de transacciones de crédito del día. & P0 & Implementado (CLIENT-01) \\
RF-15 & Consulta de totales acumulados y descarga de detalle de ventas. & P1 & Implementado (CLIENT-01) \\
RF-16 & Operaciones multicódigo para centros comerciales y comisionistas. & P1 & Implementado (CLIENT-01) \\
\end{longtable}

\subsection{Requisitos no funcionales}
La Tabla~\ref{tab:requisitos-no-funcionales} resume las restricciones técnicas y arquitectónicas.

\begin{longtable}{@{}p{1.6cm}p{7.4cm}p{2.2cm}p{3.0cm}@{}}
\caption{Requisitos no funcionales de calidad y seguridad.}\label{tab:requisitos-no-funcionales}\\
\toprule
\textbf{ID} & \textbf{Descripción no funcional} & \textbf{Categoría} & \textbf{Regla del repo} \\
\midrule
\endfirsthead
\toprule
\textbf{ID} & \textbf{Descripción no funcional} & \textbf{Categoría} & \textbf{Regla del repo} \\
\midrule
\endhead
\bottomrule
\endlastfoot
RNF-01 & Cero código inseguro (\texttt{\#![forbid(unsafe\_code)]}) en todo el workspace. & Robustez & \texttt{SAFE-01} \\
RNF-02 & Ausencia total de pánicos en producción (\texttt{unwrap}/\texttt{expect} denegados). & Estabilidad & AGENTS.md §6 \\
RNF-03 & Representación de montos como enteros en CLP (prohibido float \texttt{f32}/\texttt{f64}). & Integridad & \texttt{API-03} \\
RNF-04 & Búferes estrictamente acotados (payload máximo 4096 bytes). & Rendimiento & \texttt{SAFE-02} \\
RNF-05 & Redacción estricta de números de cuenta y vouchers en trazas y \texttt{Debug}. & Privacidad & \texttt{SEC-02}, \texttt{SEC-03} \\
RNF-06 & Determinismo absoluto en pruebas unitarias mediante tiempo inyectable. & Calidad QA & \texttt{QA-03} \\
RNF-07 & Tiempos de ejecución de la suite de pruebas inferior a 60 segundos. & Eficiencia & \texttt{QA-04} \\
RNF-08 & Dual licensing permisivo compatible con ecosistema corporativo. & Legal & \texttt{LEG-01} \\
\end{longtable}

\clearpage
\section{Arquitectura del software y diseño de componentes}

\subsection{Visión arquitectónica global}
La arquitectura de \texttt{transbank-SDK} sigue el patrón de separación de responsabilidades en capas
concéntricas y desacopladas. La Figura~\ref{fig:arch-contexto} ilustra la interacción del sistema con
los actores externos, mientras que la Figura~\ref{fig:arch-componentes} descompone los módulos internos.

\begin{figure}[htbp]
  \centering
  \includegraphics[width=0.92\textwidth]{contexto.png}
  \caption{Diagrama de contexto del sistema (C4 Nivel 1).}
  \label{fig:arch-contexto}
\end{figure}

\begin{figure}[htbp]
  \centering
  \includegraphics[width=0.92\textwidth]{componentes.png}
  \caption{Diagrama de componentes del workspace transbank-SDK (C4 Nivel 2-3).}
  \label{fig:arch-componentes}
\end{figure}

\subsection{Descomposición de capas}
\begin{enumerate}
  \item \textbf{Capa de Protocolo Sans-IO (\texttt{transbank\_pos::protocol}):} es el núcleo funcional puro.
    No realiza operaciones de lectura o escritura en archivos, sockets o puertos serie, ni consulta la hora
    del sistema operativo. Se compone de:
    \begin{itemize}
      \item \texttt{lrc}: cálculo del byte de redundancia longitudinal mediante operación XOR iterativa.
      \item \texttt{frame}: serialización de payloads delimitados por caracteres de control ASCII (STX, ETX).
      \item \texttt{decoder}: máquina de estados por streaming (\texttt{FrameDecoder}) que recibe slices de bytes
        arbitrarios y emite eventos de enlace estructurados (\texttt{LinkEvent}).
      \item \texttt{fields}: funciones auxiliares de despiece, recorte y conversión numérica tolerante.
      \item \texttt{types}: tipos de dominio encapsulados bajo el patrón \emph{newtype} (\texttt{Amount}, \texttt{Ticket}).
    \end{itemize}
  \item \textbf{Capa de Enlace (\texttt{transbank\_pos::link}):} gestiona la sincronización temporal y el protocolo
    de acuses de recibo. Coordina la espera de ACK, el envío de NAK ante tramas corruptas, los plazos de espera
    y la captura de estados intermedios (\texttt{0900}). La lógica de decisiones se centraliza en el módulo
    \texttt{exchange}, compartida entre clientes síncronos y asíncronos.
  \item \textbf{Capa de Clientes (\texttt{transbank\_pos::client}):} expone las estructuras tipadas
    \texttt{PosIntegrado} y \texttt{PosAutoservicio}. Cada método valida los argumentos de entrada, construye
    el mensaje mediante la capa de protocolo, delega el intercambio en la capa de enlace y desensambla el
    payload resultante en tipos de respuesta seguros.
  \item \textbf{Capa de Transporte (\texttt{transbank\_pos::transport}):} abstrae el medio de comunicación físico
    o virtual mediante el trait \texttt{Transport}.
\end{enumerate}

\subsection{Registro de decisiones arquitectónicas (ADR)}
La Tabla~\ref{tab:adr-resumen} resume los fundamentos técnicos de diseño formalizados en el proyecto.

\begin{longtable}{@{}p{1.4cm}p{4.0cm}p{4.2cm}p{4.8cm}@{}}
\caption{Registro de decisiones arquitectónicas (ADR).}\label{tab:adr-resumen}\\
\toprule
\textbf{ADR} & \textbf{Decisión de diseño} & \textbf{Motivación técnica} & \textbf{Consecuencia / Beneficio} \\
\midrule
\endfirsthead
\toprule
\textbf{ADR} & \textbf{Decisión de diseño} & \textbf{Motivación técnica} & \textbf{Consecuencia / Beneficio} \\
\midrule
\endhead
\bottomrule
\endlastfoot
ADR-01 & Núcleo sans-IO en módulo \texttt{protocol}. & Desacoplar algoritmos de protocolo de llamadas I/O y del reloj. & Pruebas exhaustivas deterministas y reutilización en simulador y async. \\
ADR-02 & API bloqueante síncrona por defecto. & Las operaciones en el POS duran segundos y son estrictamente secuenciales. & Simplicidad de uso directo en hilos secundarios de apps de escritorio (Tauri). \\
ADR-03 & Rechazo financiero no es \texttt{Err}. & Una tarjeta declinada es un resultado comercial normal. & El canal y el enlace permanecen intactos; no se contamina el manejo de errores I/O. \\
ADR-04 & Prohibición de reintento post-ACK. & Reenviar una venta en vuelo puede generar un doble cobro al cliente. & Fuerza al host a consultar \texttt{last\_sale()} garantizando idempotencia. \\
ADR-05 & Decodificador por máquina de estados. & Los buffers serie pueden fragmentar tramas en bytes arbitrarios. & Robusto ante trozos parciales y valores de LRC iguales a bytes de control. \\
ADR-06 & ISO-8859-1 en entrada, ASCII en salida. & Los comprobantes impresos por POS chilenos contienen tildes y caracteres locales. & Previene errores de parseo UTF-8 en respuestas; emisión limpia hacia el POS. \\
ADR-07 & Fechas como tipos locales propios. & El terminal informa fecha y hora sin información de huso horario. & Evita dependencias pesadas como \texttt{chrono} y asunciones erróneas de zona. \\
ADR-08 & Trait \texttt{Transport} con timeout. & Homogeneizar lectura serie, sockets TCP y búferes en memoria. & Permite alternar hardware real y simulador sin cambiar el enlace. \\
ADR-09 & Reloj inyectable mediante \texttt{Clock}. & Los timeouts de venta duran 150 segundos en producción. & Pruebas unitarias de timeouts sin retrasos artificiales (\texttt{sleep}). \\
ADR-10 & Simulador como crate independiente. & Facilitar pruebas end-to-end a consumidores externos del SDK. & Disponible como biblioteca y como binario de red (\texttt{tbk-pos-sim}). \\
ADR-11 & Dual licensing MIT / Apache-2.0. & Convención estándar del ecosistema Rust corporativo. & Permite integración estática libre sin restricciones virales de copyleft. \\
\end{longtable}

\clearpage
\section{Stack tecnológico detallado y justificación de componentes}

\subsection{Lenguaje de programación y compilador}
\begin{itemize}
  \item \textbf{Rust 2024 (MSRV 1.85):} el proyecto exige Rust estable con versión mínima soportada 1.85
    (\texttt{rust-version = "1.85"} en \texttt{[workspace.package]}). Rust garantiza seguridad de memoria en tiempo
    de compilación sin sobrecarga de recolector de basura (garbage collector), fundamental para bibliotecas
    de control de hardware y bajo nivel.
  \item \textbf{Resolver Cargo versión 3:} asegura resolución determinista y unificada de dependencias en
    todo el workspace multirate.
\end{itemize}

\subsection{Dependencias externas autorizadas}
La lista de dependencias está estrictamente cerrada (\texttt{.ai-context/architecture.md} §3). Cada pieza
cumple un propósito específico y verificado:
\begin{itemize}
  \item \textbf{\texttt{thiserror = "2"}:} utilizado en \texttt{transbank-pos} para generar implementaciones
    eficientes del trait \texttt{std::error::Error} con cero costo en runtime y sin allocaciones dinámicas.
  \item \textbf{\texttt{serialport = "4"}:} biblioteca multiplataforma para interactuar con puertos serie
    físicos (RS-232 y emulación USB CDC-ACM), listar puertos del sistema y consultar metadatos USB (Vendor ID,
    Product ID). Se encapsula tras la feature opcional por defecto \texttt{serial}.
  \item \textbf{\texttt{serde = "1"} y \texttt{serde\_json = "1"}:} feature opcional \texttt{serde} que habilita
    la serialización automática de respuestas y eventos, facilitando el puenteado hacia frontends en TypeScript
    (Tauri) y la salida formateada en la CLI.
  \item \textbf{\texttt{tracing = "0.1"} y \texttt{tracing-subscriber = "0.3"}:} instrumentación estructurada
    para depuración de tramas en bajo nivel. Su costo en CPU es cero cuando no hay suscriptores activos.
  \item \textbf{\texttt{clap = "4"}:} parser de argumentos de línea de comandos en modo declarativo (\texttt{derive})
    utilizado en los binarios \texttt{tbk-pos} y \texttt{tbk-pos-sim}.
  \item \textbf{\texttt{proptest = "1"}:} framework de pruebas basadas en propiedades utilizado como dev-dependency
    para fuzzing del decodificador de tramas y verificación de invariantes algebraicos.
  \item \textbf{\texttt{tokio = "1"} y \texttt{tokio-serial = "5"}:} runtime asíncrono activado exclusivamente
    bajo la feature \texttt{tokio} (incorporada en ASYNC-01) para entornos de servidor concurrentes.
\end{itemize}

\subsection{Exclusiones tecnológicas deliberadas}
Para mantener un perfil de seguridad impecable y un tamaño mínimo de binario, se han rechazado
deliberadamente las siguientes dependencias comunes:
\begin{itemize}
  \item \textbf{Sin base de datos relacional (SQL):} el SDK es una biblioteca de enlace de bajo nivel sin estado
    persistente. El historial transaccional corresponde a la base de datos del host de facturación o a la
    memoria fiscal de Transbank.
  \item \textbf{Sin \texttt{anyhow}:} el uso de errores no tipados basados en \texttt{Box<dyn Error>} destruye la
    capacidad de los consumidores de inspeccionar programáticamente las fallas. Se utiliza tipado exhaustivo.
  \item \textbf{Sin \texttt{chrono} ni \texttt{time}:} los terminales POS emiten fechas sin información de huso
    horario. Introducir librerías pesadas de calendario aumenta la superficie de ataque y el tiempo de compilación.
    Se utilizan los tipos propios \texttt{PosDate} y \texttt{PosTime}.
  \item \textbf{Sin \texttt{bytes}:} las tramas físicas tienen un tamaño máximo normativo de 4096 bytes; el uso
    de slices de memoria contigua en stack y vectores nativos de Rust evita dependencias innecesarias.
  \item \textbf{Sin variables globales (\texttt{once\_cell}, \texttt{lazy\_static}):} el SDK no mantiene estado
    global oculto. Toda la configuración y los recursos residen en instancias concretas de \texttt{Link} y
    \texttt{PosIntegrado}.
\end{itemize}

\clearpage
\section{Estructura del repositorio y desglose módulo por módulo}

\subsection{Árbol de directorios del workspace}
\begin{lstlisting}[language=bash,caption={Estructura física del repositorio transbank-SDK.}]
transbank-SDK/
├── Cargo.toml                 # Configuración de workspace Cargo y lints comunes
├── Cargo.lock                 # Bloqueo de dependencias reproducibles
├── clippy.toml                # Políticas estrictas de auditoría de código
├── rustfmt.toml               # Formato de código (edición 2024)
├── .github/workflows/ci.yml   # Pipeline de integración continua
├── crates/
│   ├── transbank-pos/         # Biblioteca principal del SDK (sans-IO + clientes)
│   │   ├── Cargo.toml         # Features: serial (default), serde, tracing, tokio
│   │   └── src/
│   │       ├── lib.rs         # Re-exports de API pública y directivas forbid
│   │       ├── error.rs       # Enumeración Error y alias Result
│   │       ├── codes.rs       # ResponseCode e IntermediateStatus
│   │       ├── exchange.rs    # Máquina de intercambio síncrona/asíncrona
│   │       ├── link.rs        # Link<T, C> y configuración LinkConfig
│   │       ├── protocol/      # lrc, frame, decoder, fields, types, function_codes
│   │       ├── message/       # integrado, autoservicio, responses
│   │       ├── client/        # PosIntegrado, PosAutoservicio
│   │       ├── transport/     # Transport trait, memory, tcp, serial
│   │       └── tokio/         # AsyncLink, AsyncPosIntegrado
│   ├── transbank-pos-sim/     # Crate simulador de terminal POS
│   │   ├── Cargo.toml
│   │   └── src/               # PosSimulator, Scenario, SimTransport, servidor TCP
│   └── transbank-pos-cli/     # Herramienta de línea de comandos
│       ├── Cargo.toml
│       └── src/main.rs        # Binario tbk-pos con subcomandos clap
└── docs/                      # Documentación para humanos, diagramas e informes
\end{lstlisting}

\subsection{Desglose de módulos en transbank-pos}
\begin{itemize}
  \item \textbf{\texttt{protocol::lrc}:} contiene la función \texttt{calculate\_lrc(data: \&[u8]) -> u8},
    que computa la suma de verificación mediante XOR acumulativo. Es utilizada por codificadores,
    decodificadores y por el simulador.
  \item \textbf{\texttt{protocol::decoder}:} implementa la estructura \texttt{FrameDecoder}. Lee flujos de
    bytes entrantes byte a byte sin descartar buffers y emite \texttt{LinkEvent::Ack}, \texttt{LinkEvent::Nak}
    o \texttt{LinkEvent::Frame(payload)}.
  \item \textbf{\texttt{protocol::types}:} expone newtypes validados en construcción. \texttt{Amount::new(monto)}
    falla inmediatamente si el valor está fuera del rango [50, 999.999.999]. \texttt{Ticket::new(str)} valida
    que los caracteres sean ASCII alfanuméricos y provee el método \texttt{padded()} para alinear a 6 posiciones.
  \item \textbf{\texttt{message::responses}:} define las estructuras \texttt{SaleResponse}, \texttt{LoadKeysResponse},
    \texttt{CloseResponse}, \texttt{TotalsResponse} y \texttt{RefundResponse}. Sobrescribe la implementación de
    \texttt{std::fmt::Debug} para asegurar que campos como el número de cuenta y el texto del voucher nunca
    se impriman en texto claro en logs de producción.
  \item \textbf{\texttt{link}:} estructura \texttt{Link<T, C>} que parametriza el transporte y el reloj.
    Implementa \texttt{send\_command} (disparo simple con espera de ACK), \texttt{request} (envío, ACK y espera
    de respuesta final con soporte de tramas 0900) y \texttt{request\_list} (acumulación de múltiples tramas
    hasta timeout de silencio).
  \item \textbf{\texttt{client::integrado}:} cliente para terminales POS Integrado. Expone métodos tipados
    como \texttt{sale}, \texttt{sale\_with\_status}, \texttt{last\_sale}, \texttt{refund}, \texttt{totals},
    \texttt{sales\_detail}, \texttt{close}, \texttt{load\_keys} y \texttt{poll}.
\end{itemize}

\clearpage
\section{Modelo de datos del protocolo y layouts de tramas}

\subsection{Estructura física de la trama}
Todas las comunicaciones en el protocolo Transbank utilizan tramas enmarcadas por caracteres de control ASCII:
\begin{center}
\texttt{<STX> + Payload + <ETX> + <LRC>}
\end{center}
\begin{itemize}
  \item \textbf{\texttt{<STX>} (0x02):} Start of Text. Delimita el inicio de una trama física.
  \item \textbf{Payload:} secuencia de campos alfanuméricos separados por el carácter pleca (\texttt{|}).
    El primer campo corresponde invariablemente al código de función de cuatro dígitos (por ejemplo, \texttt{0200}).
  \item \textbf{\texttt{<ETX>} (0x03):} End of Text. Marca el término de los datos útiles.
  \item \textbf{\texttt{<LRC>}:} Longitudinal Redundancy Check (1 byte). Checksum calculado mediante la operación
    XOR de todos los bytes comprendidos desde el primer carácter del payload hasta el byte \texttt{<ETX>} inclusive.
\end{itemize}

\subsection{Catálogo de comandos POS Integrado}
La Tabla~\ref{tab:comandos-integrado} detalla los códigos de función principales del POS Integrado.

\begin{longtable}{@{}p{1.4cm}p{3.8cm}p{3.2cm}p{6.0cm}@{}}
\caption{Comandos del protocolo POS Integrado.}\label{tab:comandos-integrado}\\
\toprule
\textbf{Código} & \textbf{Operación} & \textbf{Respuesta esperada} & \textbf{Estructura del payload de envío} \\
\midrule
\endfirsthead
\toprule
\textbf{Código} & \textbf{Operación} & \textbf{Respuesta esperada} & \textbf{Estructura del payload de envío} \\
\midrule
\endhead
\bottomrule
\endlastfoot
\texttt{0100} & Poll (sondeo) & Solo ACK (0x06) & \texttt{0100} \\
\texttt{0800} & Carga de llaves & \texttt{0810} & \texttt{0800} \\
\texttt{0200} & Venta normal & \texttt{0210} & \texttt{0200|monto|ticket|||voucher} \\
\texttt{0270} & Venta multicódigo & \texttt{0271} & \texttt{0270|monto|ticket|||voucher|codigo\_comercio} \\
\texttt{0250} & Última venta & \texttt{0260} & \texttt{0250|} \\
\texttt{0260} & Detalle de ventas & Múltiples \texttt{0261} & \texttt{0260|imprimir\_en\_pos|} \\
\texttt{1200} & Anulación de venta & \texttt{1210} & \texttt{1200|numero\_operacion} \\
\texttt{0700} & Totales del día & \texttt{0710} & \texttt{0700} \\
\texttt{0500} & Cierre de terminal & \texttt{0510} & \texttt{0500||} \\
\texttt{0300} & Modo normal & Solo ACK (0x06) & \texttt{0300} \\
\texttt{0900} & Mensaje intermedio & N/A (unidireccional) & Emitido por el POS hacia el host (\texttt{0900|codigo}) \\
\end{longtable}

\subsection{Estructura de la respuesta de venta (0210 / 0260)}
La Tabla~\ref{tab:layout-venta} documenta exhaustivamente los 21 campos que componen la respuesta de
venta (\texttt{0210}) y la consulta de última venta (\texttt{0260}).

\begin{longtable}{@{}p{0.8cm}p{4.2cm}p{1.8cm}p{1.8cm}p{4.6cm}p{1.4cm}@{}}
\caption{Layout de campos de respuesta de venta (0210 / 0260).}\label{tab:layout-venta}\\
\toprule
\textbf{Pos} & \textbf{Nombre de campo} & \textbf{Tipo} & \textbf{Longitud} & \textbf{Descripción} & \textbf{Redactado} \\
\midrule
\endfirsthead
\toprule
\textbf{Pos} & \textbf{Nombre de campo} & \textbf{Tipo} & \textbf{Longitud} & \textbf{Descripción} & \textbf{Redactado} \\
\midrule
\endhead
\bottomrule
\endlastfoot
0 & Código de función & Texto & 4 caracteres & Constante \texttt{"0210"} o \texttt{"0260"} & No \\
1 & Código de respuesta & Numérico & 2 dígitos & 00 = Aprobado, otros = Rechazo & No \\
2 & Código de comercio & Numérico & Hasta 12 & Identificador del comercio afiliado & No \\
3 & ID del terminal & Texto & Hasta 8 & Número de serie lógico del POS & No \\
4 & Ticket / Boleta & Texto & Hasta 6 & Identificador asignado por el comercio & No \\
5 & Código autorización & Texto & 6 caracteres & Código otorgado por el switch emisor & No \\
6 & Monto final cobrado & Numérico & Hasta 9 & Monto autorizado en pesos chilenos & No \\
7 & Cantidad de cuotas & Numérico & 2 dígitos & Número de cuotas pactadas & No \\
8 & Monto de la cuota & Numérico & Hasta 9 & Valor individual de cada cuota & No \\
9 & Últimos 4 dígitos & Texto & 4 caracteres & Terminación de la tarjeta bancaria & No \\
10 & Número operación & Texto & 6 dígitos & Identificador correlativo del POS & No \\
11 & Tipo de tarjeta & Texto & 2 caracteres & \texttt{CR} (Crédito), \texttt{DB} (Débito), \texttt{PR} (Prepago) & No \\
12 & Fecha contable & Texto & 8 caracteres & Formato \texttt{ddMMyyyy} asignado por Transbank & No \\
13 & Número de cuenta & Texto & Hasta 19 & Cuenta bancaria enmascarada & \textbf{Sí} \\
14 & Marca de tarjeta & Texto & 2 caracteres & \texttt{VI}, \texttt{MC}, \texttt{AX}, \texttt{DB}, \texttt{DC}, \texttt{MG} & No \\
15 & Fecha real & Texto & 8 caracteres & Fecha cronológica local (\texttt{ddMMyyyy}) & No \\
16 & Hora real & Texto & 6 caracteres & Hora cronológica local (\texttt{HHmmss}) & No \\
17 & ID del empleado & Numérico & Opcional & Identificador del cajero o turno & No \\
18 & Monto de propina & Numérico & Opcional & Propina autorizada en pesos & No \\
19 & Código de comercio prov & Numérico & Opcional & Presente en ventas multicódigo & No \\
20 & Texto del voucher & Texto & Múltiplo 40 & Líneas de comprobante fiscal impreso & \textbf{Sí} \\
\end{longtable}

\subsection{Catálogo de códigos de respuesta Transbank}
La Tabla~\ref{tab:catalogo-codigos} enumera los códigos de respuesta oficiales más frecuentes.

\begin{longtable}{@{}p{1.2cm}p{5.5cm}p{4.0cm}p{3.5cm}@{}}
\caption{Principales códigos de respuesta de Transbank.}\label{tab:catalogo-codigos}\\
\toprule
\textbf{Código} & \textbf{Mensaje oficial (Español)} & \textbf{Significado técnico} & \textbf{Acción requerida} \\
\midrule
\endfirsthead
\toprule
\textbf{Código} & \textbf{Mensaje oficial (Español)} & \textbf{Significado técnico} & \textbf{Acción requerida} \\
\midrule
\endhead
\bottomrule
\endlastfoot
\texttt{00} & Aprobado & Transacción autorizada exitosamente & Registrar pago y emitir boleta \\
\texttt{01} & Rechazado & Transacción declinada por el banco emisor & Solicitar otro medio de pago \\
\texttt{02} & Host no responde & Pérdida de enlace entre el POS y Transbank & Reintentar o verificar conexión \\
\texttt{13} & Debe ejecutar Cierre & Memoria transaccional del POS llena & Ejecutar \texttt{close()} de inmediato \\
\texttt{26} & Debe cargar llaves & Llaves criptográficas caducadas o nulas & Ejecutar \texttt{load\_keys()} \\
\texttt{67} & Cierre requerido & Switch bancario exige corte de jornada & Ejecutar \texttt{close()} \\
\texttt{74} & Error de comunicación & Falla de capa física o enlace del POS & Revisar cable serie / alimentación \\
\texttt{76} & Error de cierre & Discrepancia en la cuadratura de totales & Contactar soporte Transbank \\
\texttt{77} & Cierre no balanceado & Diferencias de conciliación financiera & Imprimir reporte y auditar caja \\
\texttt{90} & Inicialización exitosa & Código de aprobación en POS Autoservicio & Terminal listo para operar \\
\end{longtable}

\clearpage
\section{Interfaz de programación de aplicaciones (API pública)}

\subsection{Filosofía y tipos de dominio}
La API pública de \texttt{transbank-pos} sigue un diseño idiomático, declarativo y fuertemente tipado.
El uso de newtypes imposibilita errores comunes en tiempo de compilación. El siguiente fragmento
ilustra la invocación estándar de una venta en un POS Integrado.

\begin{lstlisting}[language=Rust,caption={Ejemplo completo de cobro de venta en transbank-pos.}]
use transbank_pos::{PosIntegrado, SaleRequest, IntermediateStatus};

fn main() -> Result<(), Box<dyn std::error::Error>> {
    // Apertura del puerto serie en Linux/macOS a 115.200 baudios
    let mut pos = PosIntegrado::open("/dev/ttyACM0")?;

    // Construccion y validacion rigurosa de la solicitud
    let request = SaleRequest::new(12500, "BOL101")?
        .with_voucher(true);

    println!("Iniciando cobro de $12.500...");

    // Ejecucion con callback en tiempo real para la interfaz grafica
    let response = pos.sale_with_status(&request, |status| {
        match status {
            IntermediateStatus::Code(code) => {
                println!("Estado del POS: {}", code.message());
            }
            IntermediateStatus::Unknown(raw) => {
                println!("Estado intermedio no documentado: {}", raw);
            }
        }
    })?;

    if response.is_approved() {
        println!("Venta APROBADA! Codigo Autorizacion: {}",
            response.authorization_code.unwrap_or_default());
        println!("Tarjeta: {:?} termina en {}",
            response.card_brand, response.last4_digits.unwrap_or_default());
    } else {
        println!("Venta RECHAZADA. Motivo: {}", response.response_code);
    }

    Ok(())
}
\end{lstlisting}

\subsection{Autoconexión y detección de puertos}
Para simplificar la configuración en tiendas con múltiples dispositivos periféricos (impresoras térmicas,
lectores de código de barras, balanzas), el SDK implementa una función de autodescubrimiento:

\begin{lstlisting}[language=Rust,caption={Detección y conexión automática al terminal POS.}]
use transbank_pos::serial;

// Escanea puertos serie, priorizando dispositivos USB, e intenta un comando poll
if let Some((pos, info)) = PosIntegrado::autoconnect(serial::DEFAULT_BAUD_INTEGRADO)? {
    println!("Terminal conectado exitosamente en puerto: {}", info.name);
    if let (Some(vid), Some(pid)) = (info.vid, info.pid) {
        println!("Dispositivo USB identificado: VID={:04x} PID={:04x}", vid, pid);
    }
} else {
    eprintln!("No se detecto ningun terminal POS conectado y listo.");
}
\end{lstlisting}

\clearpage
\section{Algoritmos y reglas de negocio críticas con ejemplos numéricos}

\subsection{Algoritmo de cálculo del Longitudinal Redundancy Check (LRC)}
El protocolo exige que cada trama física termine con un byte de paridad longitudinal (LRC).
La regla matemática formal define el LRC como el resultado de aplicar una operación binaria XOR
(\texttt{\^{}}) sucesiva sobre todos los bytes comprendidos entre el primer byte del payload
(inmediatamente posterior al carácter \texttt{<STX>} \texttt{0x02}) y el carácter delimitador
\texttt{<ETX>} (\texttt{0x03}) inclusive:
\begin{equation}
\text{LRC} = B_1 \oplus B_2 \oplus B_3 \oplus \dots \oplus B_n \oplus \text{0x03}
\end{equation}

\subsubsection{Ejemplo numérico paso a paso}
Consideremos la orden de venta por un monto de \$12.000 con ticket \texttt{"T00001"} y voucher habilitado:
\begin{itemize}
  \item \textbf{Payload:} \texttt{"0200|12000|T00001|||1"}
  \item \textbf{Secuencia de bytes (ASCII y Hexadecimal):}
\end{itemize}

\begin{longtable}{@{}p{1.5cm}p{1.8cm}p{2.5cm}p{8.0cm}@{}}
\caption{Traza de cálculo de LRC para payload de venta.}\label{tab:traza-lrc}\\
\toprule
\textbf{Carácter} & \textbf{Byte (Hex)} & \textbf{Acumulador XOR} & \textbf{Explicación binaria del acumulador} \\
\midrule
\endfirsthead
\toprule
\textbf{Carácter} & \textbf{Byte (Hex)} & \textbf{Acumulador XOR} & \textbf{Explicación binaria del acumulador} \\
\midrule
\endhead
\bottomrule
\endlastfoot
\texttt{'0'} & \texttt{0x30} & \texttt{0x30} & Inicialización con primer byte (\texttt{0011 0000}) \\
\texttt{'2'} & \texttt{0x32} & \texttt{0x02} & \texttt{0x30 \^{} 0x32 = 0011 0000 \^{} 0011 0010 = 0000 0010} \\
\texttt{'0'} & \texttt{0x30} & \texttt{0x32} & \texttt{0x02 \^{} 0x30 = 0011 0010} \\
\texttt{'0'} & \texttt{0x30} & \texttt{0x02} & \texttt{0x32 \^{} 0x30 = 0000 0010} \\
\texttt{'|'} & \texttt{0x7C} & \texttt{0x7E} & \texttt{0x02 \^{} 0x7C = 0111 1110} \\
\texttt{'1'} & \texttt{0x31} & \texttt{0x4F} & \texttt{0x7E \^{} 0x31 = 0100 1111} \\
\texttt{'2'} & \texttt{0x32} & \texttt{0x7D} & \texttt{0x4F \^{} 0x32 = 0111 1101} \\
\texttt{'0'} & \texttt{0x30} & \texttt{0x4D} & \texttt{0x7D \^{} 0x30 = 0100 1101} \\
\texttt{'0'} & \texttt{0x30} & \texttt{0x7D} & \texttt{0x4D \^{} 0x30 = 0111 1101} \\
\texttt{'0'} & \texttt{0x30} & \texttt{0x4D} & \texttt{0x7D \^{} 0x30 = 0100 1101} \\
\texttt{'|'} & \texttt{0x7C} & \texttt{0x31} & \texttt{0x4D \^{} 0x7C = 0011 0001} \\
\texttt{'T'} & \texttt{0x54} & \texttt{0x65} & \texttt{0x31 \^{} 0x54 = 0110 0101} \\
\texttt{'0'} & \texttt{0x30} & \texttt{0x55} & \texttt{0x65 \^{} 0x30 = 0101 0101} \\
\texttt{'0'} & \texttt{0x30} & \texttt{0x65} & \texttt{0x55 \^{} 0x30 = 0110 0101} \\
\texttt{'0'} & \texttt{0x30} & \texttt{0x55} & \texttt{0x65 \^{} 0x30 = 0101 0101} \\
\texttt{'0'} & \texttt{0x30} & \texttt{0x65} & \texttt{0x55 \^{} 0x30 = 0110 0101} \\
\texttt{'1'} & \texttt{0x31} & \texttt{0x54} & \texttt{0x65 \^{} 0x31 = 0101 0100} \\
\texttt{'|'} & \texttt{0x7C} & \texttt{0x28} & \texttt{0x54 \^{} 0x7C = 0010 1000} \\
\texttt{'|'} & \texttt{0x7C} & \texttt{0x54} & \texttt{0x28 \^{} 0x7C = 0101 0100} \\
\texttt{'|'} & \texttt{0x7C} & \texttt{0x28} & \texttt{0x54 \^{} 0x7C = 0010 1000} \\
\texttt{'1'} & \texttt{0x31} & \texttt{0x19} & \texttt{0x28 \^{} 0x31 = 0001 1001} \\
\texttt{<ETX>} & \texttt{0x03} & \textbf{\texttt{0x1A}} & \texttt{0x19 \^{} 0x03 = 0001 1010 (Checksum final: 0x1A)} \\
\end{longtable}

La trama final física transmitida al puerto serie contiene 24 bytes exactos:
\begin{center}
\texttt{[0x02, '0', '2', '0', '0', '|', '1', '2', '0', '0', '0', '|', 'T', '0', '0', '0', '0', '1', '|', '|', '|', '1', 0x03, 0x1A]}
\end{center}

\subsection{Regla de negocio crítica: ``Nunca Doble Cobro''}
La regla más delicada en terminales de cobro bancario establece que \textbf{bajo ninguna circunstancia}
se debe reintentar una venta que ya haya sido acusada de recibo (ACK) por el POS.

\subsubsection{Caso de falla operativo}
\begin{enumerate}
  \item El cajero cobra una venta de \$15.000 (boleta \texttt{"BOL101"}).
  \item El POS acusa recibo con \texttt{0x06} (ACK). En este momento, el POS ha bloqueado su pantalla,
    solicita la tarjeta del cliente y contacta al switch adquirente por su red 3G.
  \item Transcurren 150 segundos y la red adquirente sufre lentitud extrema, o bien el cable USB sufre
    una caída de tensión transitoria. El SDK agota su temporizador y emite \texttt{Err(Error::ResponseTimeout)}.
  \item \textbf{Conducta prohibida:} si el software de caja reintentara automáticamente la venta enviando
    una nueva trama \texttt{0200}, y la primera venta se hubiese aprobado silenciosamente en Transbank,
    el cliente sufriría un \textbf{doble cobro} de \$15.000 (\$30.000 en total).
  \item \textbf{Conducta correcta y forzada por el SDK:}
    \begin{itemize}
      \item El SDK entrega un error de timeout indicando que el resultado es desconocido.
      \item El host invoca inmediatamente \texttt{PosIntegrado::last\_sale()} (comando \texttt{0250}).
      \item Si la última venta en el POS tiene ticket \texttt{"BOL101"}, monto \$15.000 y código \texttt{00},
        se concluye con certeza absoluta que la venta sí se aprobó. El host guarda el voucher y entrega la boleta.
      \item Si el ticket no coincide o el código es de rechazo, el cobro original no se aplicó y es seguro
        ofrecer al cajero la posibilidad de reintentar el pago con un ticket nuevo.
    \end{itemize}
\end{enumerate}

\clearpage
\section{Seguridad de la información y protección de la privacidad}

\subsection{Modelo de amenazas STRIDE}
La Tabla~\ref{tab:stride-resumen} analiza las amenazas evaluadas sobre el canal de comunicaciones serie y red.

\begin{longtable}{@{}p{2.8cm}p{4.2cm}p{6.8cm}@{}}
\caption{Modelo de amenazas STRIDE aplicado a transbank-SDK.}\label{tab:stride-resumen}\\
\toprule
\textbf{Categoría STRIDE} & \textbf{Amenaza específica} & \textbf{Mitigación técnica en el SDK} \\
\midrule
\endfirsthead
\toprule
\textbf{Categoría STRIDE} & \textbf{Amenaza específica} & \textbf{Mitigación técnica en el SDK} \\
\midrule
\endhead
\bottomrule
\endlastfoot
\textbf{Spoofing} (Suplantación) & Inyección de tramas de aprobación falsas desde red local. & El simulador TCP escucha estrictamente en \texttt{127.0.0.1} por defecto. Se exige flag explícito para exponer interfaz de red (\texttt{SEC-04}). \\
\textbf{Tampering} (Manipulación) & Alteración de montos o tickets en tránsito en el cable serie. & Verificación estricta de redundancia longitudinal (LRC) en cada trama. Rechazo con NAK ante discrepancia de bits. \\
\textbf{Repudiation} (Repudio) & Negación de autorización de una venta realizada. & Preservación del payload original crudo (\texttt{raw}) y registro del código de autorización bancario. \\
\textbf{Information Disclosure} (Fuga) & Exposición de datos de tarjeta o vouchers en archivos de trazas. & Redacción obligatoria de número de cuenta y texto de voucher en implementaciones de \texttt{Debug} y \texttt{tracing}. \\
\textbf{Denial of Service} (DoS) & Desbordamiento de memoria por inyección de streams continuos de bytes. & Decodificador por streaming con buffers acotados (\texttt{SAFE-02}). Límite de 4096 bytes por trama. \\
\textbf{Elevation of Privilege} & Ejecución de código malicioso mediante explotación de memoria. & Prohibición absoluta de bloques \texttt{unsafe} en el código (\texttt{\#![forbid(unsafe\_code)]}). \\
\end{longtable}

\subsection{Aislamiento y redacción de datos sensibles}
En conformidad con las restricciones \texttt{SEC-02} y \texttt{SEC-03}, el SDK redacta automáticamente
toda información potencialmente sensible en sus representaciones de depuración. La estructura
\texttt{SaleResponse} muestra únicamente la longitud en caracteres del voucher y oculta la cuenta:

\begin{lstlisting}[language=Rust,caption={Comportamiento de redacción en SaleResponse::fmt.}]
// Salida redactada por defecto al imprimir con {:?} en logs de aplicacion:
SaleResponse {
    function_code: "0210",
    response_code: 00 (Aprobado),
    ticket: Some("BOL101"),
    amount: Some(15000),
    last4_digits: Some("1234"),
    account_number: Some("<redacted 19>"),
    voucher: Some("<voucher: 320 chars>"),
    ...
}
\end{lstlisting}

\clearpage
\section{Estrategia de pruebas y aseguramiento de la calidad}

\subsection{Pirámide de pruebas del workspace}
El aseguramiento de la calidad se organiza en cuatro niveles de verificación rigurosos:
\begin{enumerate}
  \item \textbf{Pruebas unitarias de algoritmos sans-IO:} verificación de vectores de prueba conocidos
    byte a byte para LRC, framing y campos de mensaje.
  \item \textbf{Pruebas basadas en propiedades (\texttt{proptest}):} generación de miles de secuencias
    arbitrarias de bytes para comprobar que el decodificador nunca entra en pánico ni cicla infinitamente.
  \item \textbf{Pruebas de integración cliente-simulador:} ejecución de intercambios transaccionales completos
    (ventas, anulaciones, cierres, timeouts, corrupción de tramas) sobre canales virtuales \texttt{SimTransport}.
  \item \textbf{Pruebas de interfaz CLI:} ejecución de subcomandos reales mediante el ejecutable compilado
    enlazado con el simulador en memoria.
\end{enumerate}

\subsection{Desglose de pruebas por suite}
La Tabla~\ref{tab:desglose-tests} resume la distribución real de las 401 pruebas ejecutadas por \texttt{cargo test}.

\begin{longtable}{@{}p{5.5cm}p{5.5cm}p{3.0cm}@{}}
\caption{Distribución de pruebas automatizadas en el workspace.}\label{tab:desglose-tests}\\
\toprule
\textbf{Suite de pruebas} & \textbf{Archivo / Crate} & \textbf{Pruebas aprobadas} \\
\midrule
\endfirsthead
\toprule
\textbf{Suite de pruebas} & \textbf{Archivo / Crate} & \textbf{Pruebas aprobadas} \\
\midrule
\endhead
\bottomrule
\endlastfoot
Unitarias de protocolo y tipos & \texttt{transbank\_pos (lib.rs)} & 75 pruebas \\
Enlace asíncrono y propiedades & \texttt{tests/async\_link.rs} & 46 pruebas \\
Mensajería de Autoservicio & \texttt{tests/autoservicio\_messages.rs} & 13 pruebas \\
Cliente Asíncrono Tokio & \texttt{tests/client\_async.rs} & 13 pruebas \\
Cliente Autoservicio & \texttt{tests/client\_autoservicio.rs} & 20 pruebas \\
Cliente Integrado & \texttt{tests/client\_integrado.rs} & 17 pruebas \\
Propiedades del Decodificador & \texttt{tests/decoder\_properties.rs} & 3 pruebas \\
Mensajería POS Integrado & \texttt{tests/integrado\_messages.rs} & 3 pruebas \\
Verificación de Licencias & \texttt{tests/license\_files.rs} & 1 prueba \\
Lógica de Enlace y ACK/NAK & \texttt{tests/link.rs} & 22 pruebas \\
Casos Borde de Enlace & \texttt{tests/link\_edge\_cases.rs} & 18 pruebas \\
Verificación basada en Modelo & \texttt{tests/model\_based.rs} & 1 prueba \\
Propiedades de Serialización & \texttt{tests/properties.rs} & 16 pruebas \\
Compatibilidad Serde JSON & \texttt{tests/serde.rs} & 6 pruebas \\
Transporte TCP de red & \texttt{tests/transport\_tcp.rs} & 5 pruebas \\
Unitarias CLI \texttt{tbk-pos} & \texttt{transbank-pos-cli (main.rs)} & 9 pruebas \\
Integración CLI & \texttt{transbank-pos-cli/tests/cli.rs} & 28 pruebas \\
CLI Autoservicio & \texttt{tests/cli\_autoservicio.rs} & 16 pruebas \\
Pruebas de Humo (Smoke) & Varios crates & 4 pruebas \\
Simulador POS y TCP Server & \texttt{transbank-pos-sim} & 56 pruebas \\
Doctests de documentación & \texttt{transbank-pos} y \texttt{sim} & 39 pruebas \\
\midrule
\textbf{Total general verificado} & \textbf{Workspace completo} & \textbf{401 pruebas (0 fallas)} \\
\bottomrule
\end{longtable}

\subsection{Comprobación de mutaciones con cargo-mutants}
Durante la fase de endurecimiento (\texttt{TEST-01}), el código de producción fue sometido a análisis de
mutaciones sintácticas con la herramienta \texttt{cargo-mutants}. Se inyectaron 654 mutaciones en los
módulos de enlace, redacción, decodificación y aprobación. \textbf{El 100\% de los mutantes viables fueron
detectados y eliminados} por la suite de pruebas (0 escapes), demostrando que ningún assert es cosmético.

\clearpage
\section{Integración continua y entrega continua (CI/CD)}

\subsection{Pipeline en GitHub Actions}
El repositorio implementa una suite de integración continua definida en \texttt{.github/workflows/ci.yml}.
El pipeline se ejecuta ante eventos de push en las ramas \texttt{main} y \texttt{develop}, así como en
disparos manuales (\texttt{workflow\_dispatch}).

\subsection{Pasos de verificación del CI}
Cada corrida del pipeline ejecuta rigurosamente la secuencia de verificación estipulada en §8 de \texttt{AGENTS.md}:
\begin{enumerate}
  \item \textbf{Verificación de formato de código:} \texttt{cargo fmt --all --check}.
  \item \textbf{Auditoría de lints completa:} \texttt{cargo clippy --workspace --all-targets --all-features -- -D warnings}.
  \item \textbf{Auditoría sin features opcionales:} \texttt{cargo clippy -p transbank-pos --no-default-features -- -D warnings}.
  \item \textbf{Ejecución de pruebas automatizadas:} \texttt{cargo test --workspace --all-features}.
  \item \textbf{Generación de documentación sin advertencias:} \texttt{RUSTDOCFLAGS="-D warnings" cargo doc --workspace --no-deps --all-features}.
  \item \textbf{Auditoría de sincronización documental:} \texttt{docs/build-docs.sh --check}.
\end{enumerate}

\section{Configuración del sistema y variables de entorno}

La Tabla~\ref{tab:variables-entorno} documenta las variables de entorno reconocidas por el repositorio.

\begin{longtable}{@{}p{4.2cm}p{3.2cm}p{6.6cm}@{}}
\caption{Variables de entorno reconocidas por el workspace.}\label{tab:variables-entorno}\\
\toprule
\textbf{Variable} & \textbf{Valor por defecto} & \textbf{Propósito y contexto} \\
\midrule
\endfirsthead
\toprule
\textbf{Variable} & \textbf{Valor por defecto} & \textbf{Propósito y contexto} \\
\midrule
\endhead
\bottomrule
\endlastfoot
\texttt{RUST\_LOG} & Vacío (desactivado) & Nivel de detalle de eventos de depuración estructurada (\texttt{error}, \texttt{warn}, \texttt{info}, \texttt{debug}, \texttt{trace}). \\
\texttt{PLANTUML} & Detección automática & Ruta explícita al binario ejecutable de PlantUML para regeneración de diagramas. \\
\texttt{TBK\_POS\_SIM\_ADDR} & \texttt{127.0.0.1:7777} & Dirección de escucha por defecto del servidor simulador TCP. \\
\end{longtable}

\clearpage
\section{Despliegue, instalación y distribución}

\subsection{Inclusión como dependencia en proyectos Cargo}
Para consumir el SDK en un software de caja o kiosco en Rust, se incorpora en el archivo \texttt{Cargo.toml}:

\begin{lstlisting}[language=bash,caption={Configuración de dependencias en proyectos consumidores.}]
[dependencies]
transbank-pos = { version = "1.0.0", features = ["serial", "serde"] }
\end{lstlisting}

\subsection{Configuración de permisos de puerto serie en el sistema operativo}
La comunicación directa con el hardware requiere permisos de acceso a los nodos de dispositivo:
\begin{itemize}
  \item \textbf{Linux (udev rules):} por defecto, las distribuciones Linux restringen el acceso a dispositivos
    serie a usuarios miembros del grupo \texttt{dialout} o \texttt{uucp}. Para otorgar acceso permanente al
    software de facturación sin privilegios de administrador, se crea una regla en \texttt{/etc/udev/rules.d/99-transbank.rules}:
\begin{lstlisting}[language=bash,caption={Regla udev para terminales Transbank en Linux.}]
# Terminal POS Transbank (ejemplo VID FTDI / Silicon Labs)
SUBSYSTEM=="tty", ATTRS{idVendor}=="0403", MODE="0660", GROUP="dialout"
\end{lstlisting}
    El usuario debe pertenecer al grupo correspondiente: \texttt{sudo usermod -aG dialout \$USER}.
  \item \textbf{macOS:} los dispositivos USB CDC-ACM y adaptadores FTDI son expuestos de forma automática bajo
    el subsistema IOKit como \texttt{/dev/cu.usbserial-*} o \texttt{/dev/cu.usbmodem*}. No se requieren reglas adicionales.
  \item \textbf{Windows:} el sistema asigna nombres correlativos de puerto COM (\texttt{COM3}, \texttt{COM4}) a
    través del controlador \texttt{usbser.sys}.
\end{itemize}

\clearpage
\section{Operación y manual de runbooks para soporte técnico}

\subsection{Rutina de apertura diaria y carga de llaves (Código 26)}
\textbf{Síntoma:} el terminal rechaza operaciones de venta con el código de error \texttt{26 (Debe cargar llaves)}.
\textbf{Causa:} expiración o revocación de las llaves de sesión criptográfica entre el terminal y Transbank.
\textbf{Solución:}
\begin{enumerate}
  \item Ejecutar desde la línea de comandos de la caja:
\begin{lstlisting}[language=bash]
tbk-pos --port /dev/ttyACM0 load-keys
\end{lstlisting}
  \item Verificar que la respuesta devuelva \texttt{Código 00 (Aprobado)}.
  \item Proceder con las operaciones habituales de venta.
\end{enumerate}

\subsection{Rutina de cierre de terminal y arqueo de caja (Códigos 13, 67, 74, 76, 77)}
\textbf{Síntoma:} el POS rechaza ventas con código \texttt{13 (Debe ejecutar Cierre)} o \texttt{67 (Cierre requerido)}.
\textbf{Causa:} la memoria interna del terminal ha alcanzado el límite máximo diario o el switch adquirente
exige el corte contable de la jornada.
\textbf{Solución:}
\begin{enumerate}
  \item Imprimir o consultar previamente los acumulados del día:
\begin{lstlisting}[language=bash]
tbk-pos --port /dev/ttyACM0 totals
\end{lstlisting}
  \item Ejecutar el comando formal de cierre de terminal:
\begin{lstlisting}[language=bash]
tbk-pos --port /dev/ttyACM0 close --yes
\end{lstlisting}
  \item El POS imprimirá el comprobante de cierre físico y reiniciará sus acumuladores contables a cero.
\end{enumerate}

\subsection{Procedimiento de recuperación ante corte de comunicación en venta}
\textbf{Síntoma:} la aplicación de caja informa \texttt{Error::ResponseTimeout} o el cable serie fue desconectado
durante el ingreso de la tarjeta.
\textbf{Solución:}
\begin{enumerate}
  \item \textbf{NO reintentar el cobro inmediatamente.}
  \item Consultar la última venta asentada en el POS:
\begin{lstlisting}[language=bash]
tbk-pos --port /dev/ttyACM0 last-sale
\end{lstlisting}
  \item Comparar el número de ticket y el monto reportado con la venta que sufrió la interrupción.
  \item Si coinciden y el código es \texttt{00}, ingresar manualmente la venta como aprobada en el sistema de caja.
  \item Si no coinciden o el código es de rechazo, proceder a cobrar nuevamente con un ticket nuevo.
\end{enumerate}

\clearpage
\section{Limitaciones conocidas y gestión de deuda técnica}

\subsection{Brechas conocidas del estado actual}
En concordancia con \texttt{.ai-context/implementation-status.md}, se registran las siguientes limitaciones:
\begin{itemize}
  \item \textbf{Validación con hardware físico pendiente (HW-01):} la totalidad de la lógica de protocolo,
    enlace, simulación y clientes ha sido validada contra el simulador y suites de pruebas exhaustivas. La
    certificación final con un terminal POS físico Verifone/Ingenico conectado a un puerto USB real está
    reservada al dueño del repositorio en la tarea \texttt{HW-01} antes de publicar en crates.io.
  \item \textbf{Preguntas abiertas de la especificación (OPEN-01 a OPEN-08):} existen detalles menores del
    firmware de Transbank pendientes de verificación experimental (por ejemplo, si los terminales emiten
    múltiples acuses de recibo en operaciones lentas o la tasa de baudios exacta en modelos específicos).
  \item \textbf{Código 90 en POS Integrado:} la versión 1.0.0 de \texttt{SaleResponse::is\_approved()}
    evalúa como aprobado el código de respuesta 90, diseñado originalmente para la inicialización de POS
    Autoservicio. Esta discrepancia histórica está documentada y su resolución formal se decidirá en el PR \#1.
\end{itemize}

\section{Roadmap de evolución del proyecto}

La Tabla~\ref{tab:roadmap-resumen} resume el plan maestro de tareas del repositorio.

\begin{longtable}{@{}p{1.8cm}p{7.2cm}p{2.2cm}p{2.8cm}@{}}
\caption{Roadmap maestro de evolución de transbank-SDK.}\label{tab:roadmap-resumen}\\
\toprule
\textbf{ID Tarea} & \textbf{Alcance y entregable} & \textbf{Prioridad} & \textbf{Estado} \\
\midrule
\endfirsthead
\toprule
\textbf{ID Tarea} & \textbf{Alcance y entregable} & \textbf{Prioridad} & \textbf{Estado} \\
\midrule
\endhead
\bottomrule
\endlastfoot
\texttt{INFRA-01} & Workspace Cargo, lints estrictos y pipeline de CI. & P0 & Completada \\
\texttt{PROTO-01} & Enmarcado de tramas, codificador, decodificador y LRC. & P0 & Completada \\
\texttt{PROTO-02} & Tipos, mensajes y códigos de respuesta de POS Integrado. & P0 & Completada \\
\texttt{LINK-01} & Trait Transport, reloj inyectable y capa de enlace Link. & P0 & Completada \\
\texttt{SIM-01} & Simulador sans-IO, inyector de escenarios y servidor TCP. & P0 & Completada \\
\texttt{CLIENT-01} & Cliente de alto nivel PosIntegrado con pruebas end-to-end. & P0 & Completada \\
\texttt{SERIAL-01} & Soporte de puerto serie físico, autoconnect y TCP. & P1 & Completada \\
\texttt{CLI-01} & Interfaz de línea de comandos \texttt{tbk-pos}. & P1 & Completada \\
\texttt{AUTO-01} & Soporte completo de POS Autoservicio (cliente + sim + CLI). & P2 & Completada \\
\texttt{ASYNC-01} & API asíncrona con runtime Tokio (\texttt{AsyncPosIntegrado}). & P2 & Completada \\
\texttt{HDOC-01} & Estructura documental humana, diagramas y README robusto. & P1 & Completada \\
\texttt{HDOC-02} & Informe técnico e informe de diagramas (PDF y DOCX). & P1 & \textbf{En curso} \\
\texttt{HDOC-03} & Informe legal chileno e informe de seguridad (STRIDE). & P1 & Planificada \\
\texttt{HDOC-04} & Guía de despliegue e informe de pruebas. & P1 & Planificada \\
\texttt{HW-01} & Validación en terreno con terminal POS físico real (Dueño). & P0 & Pendiente \\
\texttt{REL-01} & Publicación oficial en crates.io (versión 1.0.0). & P0 & Pendiente \\
\texttt{INT-01} & Integración en \texttt{rework-sistema-facturacion} (app Tauri). & P1 & Planificada \\
\texttt{FUT-01} & Investigación de soporte para POS Android y Bluetooth. & P2 & Futura \\
\end{longtable}

\clearpage
\section{Glosario de términos técnicos}

\begin{glosario}
  \termino{ACK}{Acknowledge (0x06). Byte de control que confirma la recepción correcta de una trama.}
  \termino{Adquirente}{Entidad financiera (Transbank) encargada de procesar las transacciones de pago con tarjeta en los comercios afiliados.}
  \termino{Baudios}{Unidad de medida de la velocidad de transmisión de datos por un canal serie físico (baud rate).}
  \termino{CDC-ACM}{Communications Device Class - Abstract Control Model. Estándar USB universal para emular un puerto serie sobre un enlace USB.}
  \termino{CLP}{Peso chileno. Moneda oficial de curso legal en Chile, indivisible en centavos para efectos de pagos presenciales.}
  \termino{DTE}{Documento Tributario Electrónico. Boleta o factura electrónica emitida formalmente ante el SII.}
  \termino{ETX}{End of Text (0x03). Byte de control que marca el fin de los datos útiles de una trama.}
  \termino{LRC}{Longitudinal Redundancy Check. Algoritmo de paridad bidireccional basado en XOR acumulativo.}
  \termino{NAK}{Negative Acknowledge (0x15). Byte de control que indica rechazo por checksum incorrecto.}
  \termino{PAN}{Primary Account Number. Número de 16 dígitos grabado en el plástico de la tarjeta bancaria.}
  \termino{PCI DSS}{Payment Card Industry Data Security Standard. Marco de seguridad global para protección de datos de tarjetahabientes.}
  \termino{PIN}{Personal Identification Number. Clave secreta numérica del titular de la tarjeta.}
  \termino{POS}{Point of Sale. Terminal electrónico físico para cobro presencial con tarjeta.}
  \termino{Sans-IO}{Patrón de diseño de software que aísla por completo el protocolo de red de las operaciones de entrada/salida y del reloj.}
  \termino{SII}{Servicio de Impuestos Internos de la República de Chile.}
  \termino{STX}{Start of Text (0x02). Byte de control que marca el inicio de una trama física.}
  \termino{Switch}{Plataforma informática centralizada de autorización de transacciones financieras.}
  \termino{Tauri}{Framework de desarrollo de aplicaciones de escritorio seguras utilizando frontend web y backend en Rust.}
\end{glosario}

\clearpage
\section{Anexo: Colección completa de diagramas de arquitectura}

Este anexo recopila las representaciones gráficas de alta resolución de los diez diagramas que sustentan
el diseño de \texttt{transbank-SDK}.

\begin{figure}[htbp]
  \centering
  \includegraphics[width=\textwidth]{contexto.png}
  \caption{Anexo 1: Contexto del sistema (C4 Nivel 1).}
\end{figure}

\begin{figure}[htbp]
  \centering
  \includegraphics[width=\textwidth]{componentes.png}
  \caption{Anexo 2: Componentes del workspace (C4 Nivel 2-3).}
\end{figure}

\begin{figure}[htbp]
  \centering
  \includegraphics[width=\textwidth]{casos-de-uso.png}
  \caption{Anexo 3: Casos de uso por actor.}
\end{figure}

\begin{figure}[htbp]
  \centering
  \includegraphics[width=\textwidth]{clases-dominio.png}
  \caption{Anexo 4: Clases y tipos del dominio.}
\end{figure}

\begin{figure}[htbp]
  \centering
  \includegraphics[width=\textwidth]{estados-enlace.png}
  \caption{Anexo 5: Máquina de estados del enlace.}
\end{figure}

\begin{figure}[htbp]
  \centering
  \includegraphics[width=\textwidth]{flujo-venta.png}
  \caption{Anexo 6: Flujo de proceso de venta.}
\end{figure}

\begin{figure}[htbp]
  \centering
  \includegraphics[width=\textwidth]{secuencia-venta.png}
  \caption{Anexo 7: Secuencia de venta exitosa en POS Integrado.}
\end{figure}

\begin{figure}[htbp]
  \centering
  \includegraphics[width=\textwidth]{secuencia-recuperacion-venta.png}
  \caption{Anexo 8: Secuencia de recuperación y nunca doble cobro.}
\end{figure}

\begin{figure}[htbp]
  \centering
  \includegraphics[width=\textwidth]{despliegue.png}
  \caption{Anexo 9: Despliegue de infraestructura y conectividad.}
\end{figure}

\begin{figure}[htbp]
  \centering
  \includegraphics[width=\textwidth]{ecosistema.png}
  \caption{Anexo 10: Integración en el ecosistema raysac.}
\end{figure}

\end{document}
